% GENERATED FILE. Edit canonical lessons, then run npm run generate:books.
% canonical-source-hash: fnv1a64:72143ad61a29e98a
% canonical-lessons: HI-C336-tapasya, HI-C336-sap, HI-C336-kosna, HI-C336-pachtana, HI-C336-ruthna

\chapter{The Curse}
\label{ch:hi-the-curse}
\hlchaptermodality{336}

\begin{chapteropening}
\textbf{By the end of this chapter:} \emph{I can tell how a curse was laid and what followed.}

Complete the last lesson of chapter 336: I can tell how a curse was laid and what followed.
\end{chapteropening}

\section[\texorpdfstring{\hi{तपस्या} (tapasyā)}{तपस्या (tapasyā)}]{\hi{तपस्या} (tapasyā) — penance; austerity}

\label{lesson:HI-C336-tapasya}



\begin{quote}
Before the new one: say the Hindi for immortal, then the Hindi for to swallow.
\end{quote}

\subsection*{You'll want to know: \hi{तपस्या}}
\textbf{\hi{तपस्या}} — \emph{tapasyā} — \textquotedblleft{}penance; austerity\textquotedblright{}.

\begin{quote}
\hi{ऋषि} \hi{ने} \hi{जंगल} \hi{में} \hi{बारह} \hi{साल} \hi{तपस्या} \hi{की।}
\end{quote}

\begin{quote}
\emph{ṛṣi ne jaṅgal meṁ bārah sāl tapasyā kī.}
\end{quote}

\begin{quote}
The sage did penance in the forest for twelve years.
\end{quote}

Feminine.

\begin{cousinweb}
\textbf{\hi{तपस्या}} (\emph{tapasyā}) is Sanskrit \emph{tapasyā} (Monier-Williams), from \emph{tapas}, heat and then the hard discipline of a holy life.
\end{cousinweb}

\begin{grammarlens}[title={What you've built}]
One more word for telling what happened.
\end{grammarlens}

\subsection*{Your turn}
\begin{itemize}
\raggedright
  \item \emph{Say it:} \emph{tapasyā}
  \item \emph{Say it:} \emph{tapasyā}, once more
  \item \emph{Say it:} \emph{ṛṣi ne jaṅgal meṁ bārah sāl tapasyā kī.} — The sage did penance in the forest for twelve years.
  \item \emph{Recall:} say the Hindi for immortal, then the Hindi for to swallow, then say \emph{tapasyā} again
\end{itemize}

\begin{tcolorbox}[breakable,colback=teal!4,colframe=teal!35!black,title={Before you move on}]
What does \hi{तपस्या} mean? (\textquotedblleft{}Penance; austerity\textquotedblright{}.) Say it once more.
\end{tcolorbox}
\section[\texorpdfstring{\hi{शाप} (śāp)}{शाप (śāp)}]{\hi{शाप} (śāp) — a curse}

\label{lesson:HI-C336-sap}



\begin{quote}
Before the new one: say the Hindi for to swallow, then the Hindi for penance.
\end{quote}

\subsection*{You'll want to know: \hi{शाप}}
\textbf{\hi{शाप}} — \emph{śāp} — \textquotedblleft{}a curse\textquotedblright{}.

\begin{quote}
\hi{ऋषि} \hi{ने} \hi{गुस्से} \hi{में} \hi{राजा} \hi{को} \hi{शाप} \hi{दिया।}
\end{quote}

\begin{quote}
\emph{ṛṣi ne gusse meṁ rājā ko śāp diyā.}
\end{quote}

\begin{quote}
In anger the sage laid a curse on the king.
\end{quote}

Masculine.

\begin{cousinweb}
\textbf{\hi{शाप}} (\emph{śāp}) is Sanskrit \emph{śāpa}, \textquotedblleft{}a curse\textquotedblright{} (Monier-Williams). Speech also says \textbf{\hi{श्राप}}. \textbf{\hi{शाप} \hi{देना}} is to lay a curse.
\end{cousinweb}

\begin{grammarlens}[title={What you've built}]
One more word for telling what happened.
\end{grammarlens}

\subsection*{Your turn}
\begin{itemize}
\raggedright
  \item \emph{Say it:} \emph{śāp}
  \item \emph{Say it:} \emph{śāp}, once more
  \item \emph{Say it:} \emph{ṛṣi ne gusse meṁ rājā ko śāp diyā.} — In anger the sage laid a curse on the king.
  \item \emph{Recall:} say the Hindi for to swallow, then the Hindi for penance, then say \emph{śāp} again
\end{itemize}

\begin{tcolorbox}[breakable,colback=teal!4,colframe=teal!35!black,title={Before you move on}]
What does \hi{शाप} mean? (\textquotedblleft{}A curse\textquotedblright{}.) Say it once more.
\end{tcolorbox}
\section[\texorpdfstring{\hi{कोसना} (kosnā)}{कोसना (kosnā)}]{\hi{कोसना} (kosnā) — to curse; to rail at}

\label{lesson:HI-C336-kosna}



\begin{quote}
Before the new one: say the Hindi for penance, then the Hindi for a curse.
\end{quote}

\subsection*{You'll want to know: \hi{कोसना}}
\textbf{\hi{कोसना}} — \emph{kosnā} — \textquotedblleft{}to curse; to rail at\textquotedblright{}.

\begin{quote}
\hi{बूढ़ी} \hi{औरत} \hi{चोर} \hi{को} \hi{कोस} \hi{रही} \hi{थी।}
\end{quote}

\begin{quote}
\emph{būṛhī aurat cor ko kos rahī thī.}
\end{quote}

\begin{quote}
The old woman was cursing the thief.
\end{quote}

Past, with \hi{ने} (he or she …): \hi{उसने} \hi{चोर} \hi{को} \hi{कोसा}, \emph{usne cor ko kosā}.

\begin{cousinweb}
\textbf{\hi{कोसना}} (\emph{kosnā}) goes back to Sanskrit \emph{krośati}, \textquotedblleft{}cries out\textquotedblright{} (CDIAL 3612). It is everyday cursing in words, not a sage's \textbf{\hi{शाप}}.
\end{cousinweb}

\begin{grammarlens}[title={What you've built}]
One more word for telling what happened.
\end{grammarlens}

\subsection*{Your turn}
\begin{itemize}
\raggedright
  \item \emph{Say it:} \emph{kosnā}
  \item \emph{Say it:} \emph{kosnā}, once more
  \item \emph{Say it:} \emph{būṛhī aurat cor ko kos rahī thī.} — The old woman was cursing the thief.
  \item \emph{Recall:} say the Hindi for penance, then the Hindi for a curse, then say \emph{kosnā} again
\end{itemize}

\begin{tcolorbox}[breakable,colback=teal!4,colframe=teal!35!black,title={Before you move on}]
What does \hi{कोसना} mean? (\textquotedblleft{}To curse; to rail at\textquotedblright{}.) Say it once more.
\end{tcolorbox}
\section[\texorpdfstring{\hi{पछताना} (pachtānā)}{पछताना (pachtānā)}]{\hi{पछताना} (pachtānā) — to regret; to repent}

\label{lesson:HI-C336-pachtana}



\begin{quote}
Before the new one: say the Hindi for a curse, then the Hindi for to curse.
\end{quote}

\subsection*{You'll want to know: \hi{पछताना}}
\textbf{\hi{पछताना}} — \emph{pachtānā} — \textquotedblleft{}to regret; to repent\textquotedblright{}.

\begin{quote}
\hi{शाप} \hi{के} \hi{बाद} \hi{राजा} \hi{बहुत} \hi{पछताया।}
\end{quote}

\begin{quote}
\emph{śāp ke bād rājā bahut pachtāyā.}
\end{quote}

\begin{quote}
After the curse the king bitterly regretted what he had done.
\end{quote}

Past (he …): \hi{वह} \hi{पछताया}, \emph{vah pachtāyā}.

\begin{cousinweb}
\textbf{\hi{पछताना}} (\emph{pachtānā}) is inherited from \emph{*paścōttāpa}, \textquotedblleft{}repentance\textquotedblright{}, a form CDIAL reconstructs (CDIAL 8010): a burning that comes after, from \emph{paścā}, behind, and \emph{tāpa}, heat.
\end{cousinweb}

\begin{grammarlens}[title={What you've built}]
One more word for telling what happened.
\end{grammarlens}

\subsection*{Your turn}
\begin{itemize}
\raggedright
  \item \emph{Say it:} \emph{pachtānā}
  \item \emph{Say it:} \emph{pachtānā}, once more
  \item \emph{Say it:} \emph{śāp ke bād rājā bahut pachtāyā.} — After the curse the king bitterly regretted what he had done.
  \item \emph{Recall:} say the Hindi for a curse, then the Hindi for to curse, then say \emph{pachtānā} again
\end{itemize}

\begin{tcolorbox}[breakable,colback=teal!4,colframe=teal!35!black,title={Before you move on}]
What does \hi{पछताना} mean? (\textquotedblleft{}To regret; to repent\textquotedblright{}.) Say it once more.
\end{tcolorbox}
\section[\texorpdfstring{\hi{रूठना} (rūṭhnā)}{रूठना (rūṭhnā)}]{\hi{रूठना} (rūṭhnā) — to sulk; to take offence}

\label{lesson:HI-C336-ruthna}



\begin{quote}
Before the new one: say the Hindi for to curse, then the Hindi for to regret.
\end{quote}

\subsection*{You'll want to know: \hi{रूठना}}
\textbf{\hi{रूठना}} — \emph{rūṭhnā} — \textquotedblleft{}to sulk; to take offence\textquotedblright{}.

\begin{quote}
\hi{रानी} \hi{राजा} \hi{से} \hi{रूठ} \hi{गई।}
\end{quote}

\begin{quote}
\emph{rānī rājā se rūṭh gaī.}
\end{quote}

\begin{quote}
The queen took offence at the king.
\end{quote}

Past (she …): \hi{वह} \hi{रूठ} \hi{गई}, \emph{vah rūṭh gaī}.

\begin{cousinweb}
\textbf{\hi{रूठना}} (\emph{rūṭhnā}) goes back to Sanskrit \emph{ruṣṭa}, \textquotedblleft{}angry, offended\textquotedblright{} (CDIAL 10791). It is a quiet anger that waits to be soothed with \textbf{\hi{मनाना}}.
\end{cousinweb}

\begin{grammarlens}[title={What you've built}]
One more word for telling what happened.
\end{grammarlens}

\subsection*{Your turn}
\begin{itemize}
\raggedright
  \item \emph{Say it:} \emph{rūṭhnā}
  \item \emph{Say it:} \emph{rūṭhnā}, once more
  \item \emph{Say it:} \emph{rānī rājā se rūṭh gaī.} — The queen took offence at the king.
  \item \emph{Recall:} say the Hindi for to curse, then the Hindi for to regret, then say \emph{rūṭhnā} again
\end{itemize}

\begin{tcolorbox}[breakable,colback=teal!4,colframe=teal!35!black,title={Before you move on}]
What does \hi{रूठना} mean? (\textquotedblleft{}To sulk; to take offence\textquotedblright{}.) Say it once more.
\end{tcolorbox}
